% !TeX program = xelatex
% !TeX encoding = UTF-8
% !TeX spellcheck = en_US
% !BIB program = biber
%% 
%% The above lines help editors like TeXstudio to automatically choose the right tools
%% to compile your LaTeX source file. If your tool does not support these magic comments,
%% you will need to make appropriate manual choices.
%% 
%% You can safely use "pdflatex" instead of "xelatex" if you prefer the pdfLaTeX toolchain.
%% However, pdfLaTeX will not be able to deliver the professional font experience that you
%% will get with XeLaTeX. You can also safely use "lualatex" instead of "xelatex" while
%% preserving the professional font experience if you prefer the LuaLaTeX toolchain.
%% 
%% _Important_: These magic comments should be on the first lines of your source file.
%% 
%%%%%%%%%%%%%%%%%%%%%%%%%%%%%%%%%%%%%%%%%%%%%%%%%%%%%%%%%%%%%%%%%%%%%%%%%%%%%%%%

%%%%%%%%%%%%%%%%%%%%%%%%%%%%%%%%%%%%%%%%%%%%%%%%%%%%%%%%%%%%%%%%%%%%%%%%%%%%%%%%
%% 
%%            JJJJ   K                         K   UUUU         UUUU  
%%            JJJJ   KKKK                   KKKK   UUUU         UUUU  
%%            JJJJ   KKKKKK               KKKKKK   UUUU         UUUU  
%%            JJJJ      KKKKKK         KKKKKK      UUUU         UUUU  
%%            JJJJ         KKKKKK   KKKKKK         UUUU         UUUU  
%%            JJJJ            KKKKKKKKK            UUUU         UUUU  
%%    JJ     JJJJJ               KKK               UUUUU       UUUUU  
%%  JJJJJJJJJJJJJ    KKKKKKKKKKKKKKKKKKKKKKKKKKK    UUUUUUUUUUUUUUU   
%%    JJJJJJJJJ      KKKKKKKKKKKKKKKKKKKKKKKKKKK      UUUUUUUUUUU     
%% 
%% This is an example file for using the JKU LaTeX technical report template
%% for your thesis.
%% 
%% Template created by Michael Roland (2021)
%% 
%%%%%%%%%%%%%%%%%%%%%%%%%%%%%%%%%%%%%%%%%%%%%%%%%%%%%%%%%%%%%%%%%%%%%%%%%%%%%%%%

%%%%%%%%%%%%%%%%%%%%%%%%%%%%%%%%%%%%%%%%%%%%%%%%%%%%%%%%%%%%%%%%%%%%%%%%%%%%%%%%
%% 
%% Document class: This is a koma-script book.
%% 
\documentclass[a4paper,oneside,10pt,ngerman,english]{scrbook}
%% 
%% The comma-separated list in square brackets are class options.
%% Useful options that you might want to use:
%% 
%% Paper size:
%%  * a4paper ... A4 paper size
%% 
%% Optimize for single-sided or double-sided printing:
%%  * oneside ... single-sided
%%  * twoside ... double-sided
%% 
%% Base font size:
%%  * 10pt ... 10-pt font is used for normal text
%%  * 11pt ... 11-pt font is used for normal text
%% 
%% Define document languages (the last specified language becomes the document default
%% language):
%%  * ngerman ... German
%%  * english ... English
%%  * ...
%% 
%% Alternate document classes: While the JKU report template supports the koma-script classes
%% `scrartcl', `scrreprt' and `scrbook', your should always use the book class `scrbook' for
%% your thesis.
%%  
%% _Important_: The document class should be the first line of LaTeX code in your main
%% source file. Do not place anything but comments / magic comments above that line (unless
%% you really know what you are doing).
%% 
%%%%%%%%%%%%%%%%%%%%%%%%%%%%%%%%%%%%%%%%%%%%%%%%%%%%%%%%%%%%%%%%%%%%%%%%%%%%%%%%

%%%%%%%%%%%%%%%%%%%%%%%%%%%%%%%%%%%%%%%%%%%%%%%%%%%%%%%%%%%%%%%%%%%%%%%%%%%%%%%%
%% 
%% Treat input files as UTF-8 encoded. Make sure to always load this when you use pdfLaTeX
%% so that pdfLaTeX knows how to read and interpret characters in this source file.
%% 
\usepackage[utf8]{inputenc}
%% 
%%%%%%%%%%%%%%%%%%%%%%%%%%%%%%%%%%%%%%%%%%%%%%%%%%%%%%%%%%%%%%%%%%%%%%%%%%%%%%%%

%%%%%%%%%%%%%%%%%%%%%%%%%%%%%%%%%%%%%%%%%%%%%%%%%%%%%%%%%%%%%%%%%%%%%%%%%%%%%%%%
%% 
%% Use the JKU LaTeX technical report template for this document.
%% 
\usepackage[mathesis,fancyfonts]{jkureport}
%% 
%% The comma-separated list in square brackets are theme options. Useful options that you
%% might want to use:
%% 
%% Document type:
%%  * phdthesis     ... PhD thesis.
%%  * mathesis      ... Master's thesis.
%%  * diplomathesis ... Diploma thesis.
%%  * bathesis      ... Bachelor's thesis.
%%  * seminarreport ... Seminar report.
%%  * techreport    ... Technical report.
%% 
%% Color scheme selection options:
%%  * JKU  ... Use JKU (gray) color scheme (this is the default if no scheme is selected).
%%  * BUS  ... Use Business School color scheme.
%%  * LIT  ... Use Linz Institute of Technology color scheme.
%%  * MED  ... Use MED faculty color scheme.
%%  * RE   ... Use RE faculty color scheme.
%%  * SOE  ... Use School of Education color scheme.
%%  * SOWI ... Use SOWI faculty color scheme.
%%  * TNF  ... Use TNF faculty color scheme.
%% 
%% Space-efficient monospace font options (requires XeTeX/LuaTeX):
%%  * compactmono   ... Use condensed fixed-width font everywhere.
%%  * nocompactverb ... Do not use condensed fixed-width font for verbatim and listings.
%% 
%% Style-breaking options:
%%  * noimprint      ... Do not insert imprint on title pages.
%%  * nojkulogo      ... Do not insert JKU & K logos on title pages.
%%  * capstitle      ... Set document title in capital letters.
%%  * nofancyfonts   ... Do not use custom TTF fonts with XeTeX/LuaTeX / supress pdfLaTeX warning.
%%  * equalmargins   ... Decrease the outer page margin to have both page margins of equal size
%%                       (the additional outer margin is intentional and to be used for
%%                       anotations; equalmargins also causes the text width to be
%%                       significantly larger than optimal for reading).
%% 
%% Experimental options:
%%  * mathastext ... Use standard document fonts (enhanced with symbols from Fira Math font
%%                   when using XeTeX/LuaTeX) in math mode.
%% 
%% Advanced options:
%%  * legacymode={<mode>} ... Activate a legacy mode to account for previous formal requirements,
%%                            styles, and/or features. May be specified mutliple times to activate
%%                            multiple legacy modes. The following <mode>s are supported:
%%      - phd2021             ... Adapt cover sheet to formal requirements for unstructured
%%                                doctoral programs before 2021W (also accounts for change from
%%                                "Beurteiler*in" to "Betreuer*in" in 2024W versions of the German
%%                                thesis cover sheets)
%%  * noautopdfinfo       ... Do not automatically try to add pdfinfo with hyperref from document
%%                            metadata fields.
%%  * logopath={<path>}   ... Set the path where the theme can find its own logo resources. This
%%                            should typically be a relative path and the default is `./logos'.
%%  * fontpath={<path>}   ... Set the path where the theme can find its own font resources. This
%%                            should typically be a relative path and the default is `./fonts'.
%% 
%% Hint: Boolean options can be used in the forms `option' or `option=true' the enable the
%% option and `nooption' or `option=false' to disable the option.
%% 
%%%%%%%%%%%%%%%%%%%%%%%%%%%%%%%%%%%%%%%%%%%%%%%%%%%%%%%%%%%%%%%%%%%%%%%%%%%%%%%%

%%%%%%%%%%%%%%%%%%%%%%%%%%%%%%%%%%%%%%%%%%%%%%%%%%%%%%%%%%%%%%%%%%%%%%%%%%%%%%%%
%% 
%% This is the place where you can load additional packages. If you want to load
%% a package `biblatex', you would use the command `\usepackage{biblatex}'.
%% 

\usepackage{csquotes}
\usepackage[backend=biber,citestyle=numeric,sortcites=true,maxcitenames=2,style=ACM-Reference-Format]{biblatex}
\setcounter{biburlnumpenalty}{100} %% reducing biburl* penalties typically improves URL placement in bibliography
\setcounter{biburllcpenalty}{100}
\setcounter{biburlucpenalty}{100}
\usepackage{todonotes}
\usepackage{import}
\usepackage{amsfonts}
\usepackage{subfigure}
%\usepackage{acronym}

%% 
%%%%%%%%%%%%%%%%%%%%%%%%%%%%%%%%%%%%%%%%%%%%%%%%%%%%%%%%%%%%%%%%%%%%%%%%%%%%%%%%

%%%%%%%%%%%%%%%%%%%%%%%%%%%%%%%%%%%%%%%%%%%%%%%%%%%%%%%%%%%%%%%%%%%%%%%%%%%%%%%%
%% 
%% Bibliography data files.
%% 

\addbibresource{references.bib}

%% 
%%%%%%%%%%%%%%%%%%%%%%%%%%%%%%%%%%%%%%%%%%%%%%%%%%%%%%%%%%%%%%%%%%%%%%%%%%%%%%%%

\begin{document}
%% Begin with the frontmatter (among other things, this switches to roman page numbers)
\frontmatter

%%%%%%%%%%%%%%%%%%%%%%%%%%%%%%%%%%%%%%%%%%%%%%%%%%%%%%%%%%%%%%%%%%%%%%%%%%%%%%%%
%% 
%% Thesis information and title page
%% 

%% Command \title{title}: sets the title of your thesis
\title{Token-Level Relevance Weighting for Uncertainty Quantification in Large Language Model Question Answering}

%% Command \titleshort{short title}: sets an abbreviated version of the thesis title for page heads
%\titleshort{Optional space for your abbreviated title}

%% Command \subtitle{subtitle}: sets the subtitle for seminar/technical reports (not used for theses)
%\subtitle{Seminar Report}

%% Command \author{name}: sets the author's name; use \prefix{} and \suffix{} to add academic titles
%%   and suffixes, use \matno{} to add the immatriculation number
\author{Nikolas Boldis \suffix{Bc.}\matno{11706481}}

%% Command \supervisor[number,gender]{name}: sets the name of the supervisor (where number optionally
%%   defines the rank of the supervisor (1-3) and gender specifies if the supervisor is male (m) or
%%   female (f/w) to adapt gender-specific terms in German)
\supervisor{\prefix{Prof. Dr.} Firstname~Lastname}

%% Command \assistantsupervisor{name}: sets the name of the assistant supervisor(s)
\assistantsupervisor{\prefix{Dr.} Firstname~Lastname \suffix{MSc}}

%% Command \degree{degree}{degree program}: sets the degree and degree program name
\degree{Master of Science}{Artificial Intelligence}
%% Please make sure to insert the _correct_ degree and degree program names for your studies before
%% submission. Examples for degrees at JKU:
%%  * Diplom-Ingenieurin / Diplom-Ingenieur
%%  * Doctor medicinae universae
%%  * Master of Arts
%%  * Master of Laws
%%  * Master of legal and business aspects in technics
%%  * Master of Science
%%  * Master of Science in Digital Business Management
%%  * Master of Science Global Business
%%  * Master of Social Sciences
%%  * Doktorin / Doktor der Geistes- und Kulturwissenschaften
%%  * Doktorin / Doktor der Medizinischen Wissenschaften
%%  * Doktorin / Doktor der Naturwissenschaften
%%  * Doktorin / Doktor der Rechtswissenschaften
%%  * Doktorin / Doktor der Sozial- und Wirtschaftswissenschaften
%%  * Doktorin / Doktor der technischen Wissenschaften
%%  * Doctor of Philosophy

%% Command \submissiondepartment{institute or department}: set the department that the thesis is submitted at
\submissiondepartment{Institute for Machine Learning}

%% Command \date{YYYY-MM-DD}: set the day of submission (defaults to today)
%\date{2020-04-09}

%% Command \keywords{text}: set the document keywords
%\keywords{Space for your comma-separated keywords}


%% Finally, print the title page using the above information:
\maketitle
%% 
%%%%%%%%%%%%%%%%%%%%%%%%%%%%%%%%%%%%%%%%%%%%%%%%%%%%%%%%%%%%%%%%%%%%%%%%%%%%%%%%

%%%%%%%%%%%%%%%%%%%%%%%%%%%%%%%%%%%%%%%%%%%%%%%%%%%%%%%%%%%%%%%%%%%%%%%%%%%%%%%%
%% 
%% Include your abstract into the frontmatter
%% 

\import{./}{00-abstract}

%% 
%%%%%%%%%%%%%%%%%%%%%%%%%%%%%%%%%%%%%%%%%%%%%%%%%%%%%%%%%%%%%%%%%%%%%%%%%%%%%%%%

%%%%%%%%%%%%%%%%%%%%%%%%%%%%%%%%%%%%%%%%%%%%%%%%%%%%%%%%%%%%%%%%%%%%%%%%%%%%%%%%
%% 
%% Include your acknowledgements into the frontmatter (optional and typically only
%% used to acknowledge external funding, discuss with your supvervisor if unsure)
%% 

%\import{./}{acknowledgements}

%% 
%%%%%%%%%%%%%%%%%%%%%%%%%%%%%%%%%%%%%%%%%%%%%%%%%%%%%%%%%%%%%%%%%%%%%%%%%%%%%%%%

%%%%%%%%%%%%%%%%%%%%%%%%%%%%%%%%%%%%%%%%%%%%%%%%%%%%%%%%%%%%%%%%%%%%%%%%%%%%%%%%
%% 
%% Add a table of contents (and optionally various other lists)
%% 

%% Make sure to start the table of contents on a new odd page (odd is only relevant in twoside layout)
\cleardoubleoddpage
%% Print the table of contents
\tableofcontents

%% Make sure to start the list of tables on a new odd page (odd is only relevant in twoside layout)
%\cleardoubleoddpage
%% Print the list of tables (optional and often not necessary)
%\listoftables

%% Make sure to start the list of figures on a new odd page (odd is only relevant in twoside layout)
%\cleardoubleoddpage
%% Print the list of figures (optional and often not necessary)
%\listoffigures

%% Make sure to start the list of acronyms on a new odd page (odd is only relevant in twoside layout)
%\cleardoubleoddpage
%% Include list of acronyms (optional and often not necessary)
%\import{./}{acronyms}

%% 
%%%%%%%%%%%%%%%%%%%%%%%%%%%%%%%%%%%%%%%%%%%%%%%%%%%%%%%%%%%%%%%%%%%%%%%%%%%%%%%%

%% Begin with the mainmatter (among other things, this switches to arabic page numbers)
\mainmatter

%%%%%%%%%%%%%%%%%%%%%%%%%%%%%%%%%%%%%%%%%%%%%%%%%%%%%%%%%%%%%%%%%%%%%%%%%%%%%%%%
%% 
%% Include your chapters
%% 


\chapter{Introduction}
\label{ch:intro}
% TODO (chapter scope): ~3-5 pages. Funnel structure: broad (LLMs in
% high-stakes use) -> narrow (UQ matters) -> specific (token aggregation
% is underspecified) -> this thesis. End the chapter with the outline.
% Do NOT cover method details, related work, or results here.
%
% NB on citations: bibkeys below are PLACEHOLDERS. Please add entries
% to references.bib before first compile -- confirm the key spellings
% match what you put in the .bib.

\section{Motivation}
\label{sec:motivation}
Large language models (LLMs) are increasingly deployed in settings where a
confident wrong answer is more harmful than an admitted absence of knowledge:
clinical decision support, legal drafting, scientific literature review, and
software engineering copilots. In all of these settings, the consumer of the
generated text treats the output as evidence and acts on it. When the model is
correct, fluency is a feature; when the model is wrong but fluent, fluency is
an attack surface. Reliable \emph{uncertainty quantification} (UQ) is therefore
a prerequisite rather than an accessory for safe LLM deployment
\cite{amodei2016concreteproblems,hendrycks2021unsolved,farquhar2024detecting}.

UQ for LLMs is hard for reasons that do not arise in standard classification
or regression. Outputs are variable-length token sequences, many of which are
semantically equivalent \cite{kuhn2023semantic}. The model samples
autoregressively, so the same prompt can produce lexically different but
semantically identical answers, and the probability the model assigns to its
own output is defined over tokens rather than meanings. Consequently, a raw
per-token probability signal captures both genuine epistemic uncertainty about
the answer and superficial lexical variation that the application does not
care about \cite{hullermeier2021aleatoric}.

Two broad families of methods dominate the current literature. \emph{Sampling-based}
methods, such as semantic entropy \cite{kuhn2023semantic} and Shifting
Attention to Relevance (SAR) \cite{duan2024sar}, generate multiple output
sequences and analyse their agreement. They are theoretically attractive but
require many forward passes per query and, in the case of semantic entropy, an
auxiliary natural-language inference model to detect semantic equivalence.
\emph{Single-pass} methods, of which the most recent is the generative
negative log-likelihood (G-NLL) proposed by Aichberger et
al.~\cite{aichberger2024rethinking}, score a single greedy decoding by summing
its per-token log-likelihoods. G-NLL is competitive with sampling-based
methods while requiring only the forward passes that were already needed to
produce the answer, which makes it the pragmatic choice for deployment at
scale.

The efficiency of G-NLL comes at a conceptual cost. G-NLL aggregates per-token
log-likelihoods uniformly, implicitly asserting that every generated token
contributes equally to the answer's uncertainty. This assumption is at odds
with the central observation of SAR \cite{duan2024sar}, namely that natural
language exhibits substantial \emph{linguistic redundancy}: a handful of
content-bearing tokens carry the answer's meaning, while function words and
boilerplate contribute little to its correctness. SAR operationalises this
observation by weighting token log-likelihoods by a relevance score obtained
through multi-sample token ablation. G-NLL ignores it. This thesis is
concerned with the question of whether the relevance-weighting idea of SAR
can be imported into the single-pass regime of G-NLL---that is, whether we
can preserve the computational profile of G-NLL while correcting its uniform
aggregation.


\section{Problem Statement}
\label{sec:problem}
% TODO: ~0.5 page. Focus: state the GAP precisely, not the solution.
Given a prompt $x$ and a greedy decoding $y = (y_1, \ldots, y_T)$ produced by
an LLM, the goal of sequence-level UQ for selective question answering is to
assign $y$ a scalar score $u(y \mid x) \in \mathbb{R}$ such that the ranking of
scores separates correct from incorrect answers. The state-of-the-art
single-pass estimator, G-NLL, defines $u$ as an unweighted sum of per-token
negative log-likelihoods \cite{aichberger2024rethinking}. This estimator is
cheap, theoretically grounded, and empirically strong, but by construction it
treats every token position as equally informative about answer correctness.

The evidence against uniform aggregation is twofold. First, the linguistic
argument made by Duan et al.~\cite{duan2024sar}: content-bearing tokens and
function words are not interchangeable carriers of meaning, so an aggregator
that treats them as such discards signal that correlates with correctness.
Second, the empirical analysis reported in
Chapter~\ref{ch:token-analysis} of this thesis, in which per-token
negative log-likelihoods are shown to vary over more than an order of
magnitude within a single answer and to differ systematically between correct
and incorrect generations.

SAR addresses uniform aggregation directly but at a cost that negates the
efficiency rationale of G-NLL: it requires $K$ stochastic samples per query
and a per-token ablation pass per sample, yielding roughly $K{\cdot}(T{+}1)$
forward passes for a single uncertainty score \cite{duan2024sar}. To date, no
single-pass estimator has combined the token-level relevance weighting of SAR
with the greedy-decoding economics of G-NLL. Closing this gap is the problem
this thesis addresses.


\section{Research Questions}
\label{sec:rqs}
% TODO: ~0.25-0.5 page. Each RQ must be measurable by ONE chapter/section:
%   RQ1 -> answered in Ch.~\ref{ch:token-analysis} (empirical NLL heterogeneity)
%   RQ2 -> answered in Ch.~\ref{ch:rwgnll} + \S6.1 (RW-G-NLL vs G-NLL AUROC)
%   RQ3 -> answered in \S6.2 (vs position-based schemes)
%   RQ4 -> answered in \S6.4 (1B / 8B / 70B scaling)
We structure the thesis around four research questions, each answered by a
dedicated chapter or results section.
\begin{description}
    \item[RQ1] Do tokens contribute unequally to the uncertainty of a greedy
    LLM answer, as measured by the within-answer distribution of per-token
    negative log-likelihoods, and does this distribution differ between
    correct and incorrect answers?
    \item[RQ2] Can a single-pass, relevance-weighted aggregator of per-token
    log-likelihoods improve selective-QA AUROC over G-NLL without additional
    sampling?
    \item[RQ3] How does such a relevance-weighted aggregator compare to
    purely positional weighting schemes that require no semantic signal?
    \item[RQ4] Do the empirical findings for RQ2--RQ3 generalise across
    model families and across model scales from 1B to 70B parameters?
\end{description}


\section{Hypothesis}
\label{sec:hypothesis}
% TODO: ~0.25 page. ONE falsifiable sentence, then one paragraph of rationale.
\paragraph{H1.}
Weighting per-token negative log-likelihoods by a semantic relevance score
derived from single-shot token ablation yields strictly higher AUROC for
selective question answering than unweighted G-NLL, across model families and
across model scales from 1B to 70B parameters, on TriviaQA.

\noindent
The hypothesis is motivated by two observations established in the
literature. First, per-token probabilities reflect both meaning-preserving
lexical variation and genuine answer uncertainty, and the former does not
correlate with correctness \cite{kuhn2023semantic}. Second, token relevance
estimated via ablation is a useful correction signal, as demonstrated by
SAR in the multi-sample regime \cite{duan2024sar}. If these two observations
transfer to the single-sample regime, a relevance-weighted variant of G-NLL
should dominate the uniform-weighted baseline. Mechanism, estimator, and
computational profile are deferred to Chapter~\ref{ch:rwgnll}.


\section{Contributions}
\label{sec:contributions}
This thesis makes the following contributions.
\begin{description}
    \item[C1] An empirical characterisation of per-token negative
    log-likelihood heterogeneity under greedy decoding, showing that token
    NLL varies substantially within answers and differs systematically
    between correct and incorrect generations (Chapter~\ref{ch:token-analysis}).
    \item[C2] \textsc{RW-G-NLL}, a single-pass, sampling-free,
    relevance-weighted aggregator of per-token log-likelihoods that preserves
    the computational profile of G-NLL while correcting its uniform
    aggregation (Chapter~\ref{ch:rwgnll}).
    \item[C3] A systematic evaluation of RW-G-NLL against G-NLL,
    length-normalised NLL, positional weighting schemes, SAR, and semantic
    entropy, across model families and model scales from 1B to 70B
    parameters on TriviaQA, with secondary validation on CoQA
    (Chapter~\ref{ch:results}).
    \item[C4] A reproducible open-source pipeline, including configuration,
    judge protocol, and per-model run scripts, covering all experiments
    reported in this thesis (Appendix~\ref{app:impl}).
\end{description}


\section{Thesis Outline}
\label{sec:outline}
% TODO: ~0.25 page. Roadmap only.
Chapter~\ref{ch:background} reviews the background required for the rest of
the thesis: the Transformer architecture, autoregressive generation, the
aleatoric--epistemic distinction as it applies to language generation, and the
three families of LLM UQ methods against which RW-G-NLL will be compared.
Chapter~\ref{ch:token-analysis} reports an exploratory empirical analysis of
per-token NLL behaviour under greedy decoding, establishing the heterogeneity
that motivates a non-uniform aggregator. Chapter~\ref{ch:rwgnll} formalises
RW-G-NLL, specifies its token-ablation relevance signal, and analyses its
computational profile relative to SAR. Chapter~\ref{ch:exp-design} describes
the full experimental design used to answer RQ1--RQ4, including the model
grid, datasets, baselines, and evaluation protocol.
Chapter~\ref{ch:results} reports the main results. Chapter~\ref{ch:discussion}
interprets those results, returns to RQ1--RQ4, and discusses limitations and
future work.

\chapter{Background}
\label{ch:background}
% TODO (chapter scope): ~7-10 pages. Goal: give an AI-MSc reader exactly
% the minimum they need to read Ch.~3-6. Be ruthlessly selective. For
% every topic ask: "does RW-G-NLL depend on this?" If no, cut it or push
% it to the appendix. Cite -- do not re-derive -- well-known results.
% Budget split (approx): \S2.1=1p, \S2.2=1.5p, \S2.3=1.5p, \S2.4=3p, \S2.5=0.5p.
%
% Suggested narrative arc:
%   \S2.1 How LLMs produce tokens (Transformer + autoregression) ->
%   \S2.2 What a probability over a token is and how G-NLL sums them ->
%   \S2.3 Aleatoric vs epistemic in the LLM setting ->
%   \S2.4 How the field currently quantifies uncertainty (G-NLL, PPL,
%         Semantic Entropy, SAR) ->
%   \S2.5 How we measure "good UQ" (AUROC for selective QA).
\section{Transformer Architecture}
% TODO: ~1 page. Focus: only what is needed for the reader to
% understand WHY a per-token probability exists. Keep the attention
% equation, one sentence on multi-head, one on residuals + LN, one on
% the final softmax head. Cite Vaswani et al. (2017). Do NOT re-derive
% backprop, explain training, or survey architecture variants.
% Transformers & LLMs (brief overview)
Start with motivation:
The Transformer architecture (Vaswani et al., 2017) revolutionized NLP by replacing recurrent structures with self-attention mechanisms, enabling parallel processing and better long-range dependency modeling.

Explain key concepts:
Self-attention mechanism (with equation)
\begin{equation}
\mathrm{Attention}(Q,K,V) = \mathrm{softmax}\left(\frac{QK^{\mathsf{T}}}{\sqrt{d_k}}\right)V
\end{equation}
Multi-head attention
Positional encoding
Feed-forward networks

Focus on what's relevant to THE work:
Crucially for our work, Transformers generate text autoregressively, computing $P(y_t \mid y_1,\dots,y_{t-1}, x)$ for each token. This token-level probability is the foundation for uncertainty quantification.
\section{Autoregressive Generation}
% TODO: ~1.5 pages. This is the technical spine the rest of the thesis
% stands on. Must define cleanly:
%   (i)   P(y_t | y_{<t}, x)
%   (ii)  greedy decoding vs stochastic sampling (temperature, top-k/p)
%   (iii) the sequence log-likelihood and its negation (G-NLL)
%   (iv)  length normalisation (mean-NLL / perplexity)
% Fix NOTATION here once and reuse it everywhere. Every later chapter
% should reference these equations by number rather than redefining.
Explain the generation process (pseudocode):
\begin{verbatim}
for t in range(max_length):
    logits = model(input_ids)
    probs = softmax(logits)
    next_token = sample(probs)  # or argmax for greedy
    input_ids.append(next_token)
\end{verbatim}

Connect to uncertainty:
Each token's log-probability $\log P(y_t\mid\text{context})$ provides a measure of the model's confidence. Aggregating these across the sequence yields the generative negative log-likelihood (G-NLL):
\begin{equation}
\mathrm{G\mbox{-}NLL}(y\mid x) = -\sum_{t=1}^{T} \log P\left(y_t\mid y_{<t}, x\right).
\end{equation}

\section{Fundamentals of Predictive Uncertainty}
% TODO: ~1.5 pages total for this section. Keep it ML-generic first,
% then specialise to LLMs.

\subsection{Aleatoric vs. Epistemic Uncertainty}
% TODO: ~0.75 page. Focus: the classical definition (Hora 1996 / Der
% Kiureghian & Ditlevsen 2009 / Hullermeier & Waegeman 2021) and how
% it maps onto LLM generation. Make clear that per-token log-probs
% mostly capture aleatoric uncertainty under a fixed model, and that
% epistemic uncertainty requires sampling/ensembling (link forward to
% Semantic Entropy in \S2.4.2). This distinction motivates why RW-G-NLL
% is a CHEAP aleatoric-style signal, not a full UQ solution.

\subsection{Uncertainty in the Context of Language Generation}
% TODO: ~0.75 page. Focus: what makes LLM UQ hard compared to classification.
% Bullets to cover: variable-length outputs, semantic equivalence
% (``Paris'' vs ``the capital of France''), exposure bias, and the
% open-endedness problem (Kuhn et al. 2023). End with a short list of
% what a "good" UQ signal must satisfy (calibration, selective QA usable,
% cheap at inference). This list is the target RW-G-NLL will be judged on.

\section{Uncertainty Quantification Methods}
% TODO (section scope): ~3 pages total. Organise by COST tier, not by year:
%   \S2.4.1 single-pass (G-NLL, PPL, length-normalised NLL) -- CHEAP, your
%           family.
%   \S2.4.2 sampling-based (Semantic Entropy, Monte Carlo predictive
%           entropy) -- EXPENSIVE, the strong baseline.
%   \S2.4.3 token-importance (SAR) -- HYBRID, your direct inspiration.
% For each: (1) the formula, (2) cost in forward passes, (3) AUROC-range
% reported in the original paper, (4) one-line strength/weakness.
\subsection{Sequence-Level Methods}
G-NLL (Your baseline):
\begin{equation}
\mathrm{G\mbox{-}NLL}(y\mid x) = -\sum_{t=1}^{T} \log P\left(y_t\mid y_{<t}, x\right).
\end{equation}
Properties:
\begin{itemize}
\item Simple to compute (no additional inference).
\item Treats all tokens equally.
\item Baseline in multiple studies [cite].
\end{itemize}
Perplexity:
\begin{equation}
\mathrm{PPL} = \exp\left(\frac{\mathrm{G\mbox{-}NLL}(y\mid x)}{T}\right),
\end{equation}
where $T$ is the number of generated tokens.


\subsection{Sampling-Based Methods}
% TODO: focus on Semantic Entropy (Kuhn et al. 2023) as the primary
% sampling-based baseline you will compare against. Explain bidirectional
% entailment clustering in 1 paragraph; do not re-prove their theorems.
% Mention predictive entropy / lexical similarity briefly as historical
% context. Be explicit about cost: SE requires ~10 samples + an NLI model
% per query -- this sets up the efficiency argument for RW-G-NLL.
Semantic Entropy (Kuhn et al., 2023):
Rather than measuring syntactic variation, semantic entropy clusters semantically equivalent answers using entailment models, then computes entropy over semantic clusters.
\begin{equation}
\mathrm{SE} = -\sum_{c} P(c)\,\log P(c),
\end{equation}
where $c$ denotes semantic clusters.
Key insight to mention:
Semantic Entropy showed that LLMs can know what they don't know, achieving AUROC > 0.7 for detecting incorrect answers on various QA tasks.
\subsection{Token Importance Methods}
% TODO: this subsection is load-bearing because RW-G-NLL is framed
% against SAR. Cover: (a) the relevance-weight definition via token
% ablation + sentence similarity, (b) the fact that SAR averages over
% MULTIPLE stochastic samples, (c) its reported AUROC ballpark. End
% the subsection with one sentence of the form "RW-G-NLL isolates SAR's
% relevance-weighting idea and applies it to a SINGLE greedy generation,
% which we formalise in Ch.~\ref{ch:rwgnll}." -- no method details here.
SAR - Shifting Attention to Relevance (The inspiration):
This is CRITICAL to cite and explain well since RW-G-NLL builds on this.
Dhuliawala et al. (2023) introduced SAR, which weights token log-likelihoods by their semantic relevance. For each token, they measure relevance via ablation: removing the token and computing semantic similarity between the original and ablated text.
\begin{equation}
\mathrm{SAR}(y\mid x) = \sum_{t=1}^{T} w_t\,\bigl(-\log P(y_t\mid y_{<t}, x)\bigr),
\qquad w_t = 1 - \mathrm{sim}\bigl(y, y_{\setminus t}\bigr).
\end{equation}
Difference from G-NLL:
\begin{itemize}
\item Requires multiple stochastic samples (expensive).
\item Computes relevance weights via ablation.
\item Normalizes by the sum of weights.
\end{itemize}
Our contribution vs SAR:
While SAR requires multiple temperature-based samples (typically 10-20), our RW-G-NLL applies the same relevance weighting principle to single greedy generations, making it computationally efficient for deployment.
% SAR (relevance weighting inspiration)
\section{Evaluation Metrics}
% TODO: ~0.5 page. Focus: justify AUROC as the PRIMARY metric for selective
% QA -- treat UQ as a binary classifier of "answer correct vs incorrect"
% using the uncertainty score as logit. Mention (briefly) why you did NOT
% use: ECE (requires calibration, not just ranking), accuracy@coverage
% curves (secondary, in appendix), Brier (needs probabilistic targets).
% Define correctness operationally via LLM-as-judge and forward-ref \S5.4
% for the exact judge protocol.
We use Area Under the ROC Curve to measure how well uncertainty scores separate correct from incorrect answers. An AUROC of 0.5 indicates random performance, while 1.0 indicates perfect separation.

\chapter{Empirical Analysis of Token-Level NLL}
\label{ch:token-analysis}
% TODO (chapter scope): ~3-5 pages. Exploratory study only; the full
% multi-model grid is in Ch.~\ref{ch:results}. Budget split: \S3.1=0.5p,
% \S3.2=2p (three figures), \S3.3=0.5p, \S3.4=0.25p.
%
% No \section{Motivation} here on purpose: motivation lives in \S1.1--\S1.2.
% Open the chapter with 2-3 sentences that cross-reference the intro and
% state the narrow question this chapter answers. The paragraph below is
% your current draft text; TODO: tighten to a cross-reference such as
% "As argued in \S\ref{...}, G-NLL aggregates token log-probabilities
% uniformly. This chapter asks whether, empirically, tokens differ enough
% in their per-token NLL to justify a non-uniform scheme."
Before proposing a weighting scheme, we must first establish whether tokens indeed contribute unequally to answer quality. We conduct an exploratory analysis of token-level negative log-likelihood patterns.
\section{Experimental Setup}
% TODO: ~0.5 page. Give EVERY knob needed to reproduce the figures
% in \S3.2: model checkpoint (HF id + revision), dataset split + size,
% decoding args (temperature=0, max_new_tokens, stop sequences),
% tokenizer, prompt template, judge model + prompt, hardware.
% Note: this is a SMALL-scale exploratory study -- make clear that the
% full multi-model evaluation lives in Ch.~\ref{ch:exp-design}. This
% avoids the reader thinking your main results rest on 1B/400 samples.
Model: Llama-3.2-1B (chosen for computational efficiency)
Dataset: TriviaQA validation set (400 samples)
Generation: Greedy decoding (temperature=0.0)
Evaluation: LLM-as-judge (Llama-3.1-70B) for binary correctness

For each generated answer, we extract:
\begin{itemize}
\item Token strings: $y_t$
\item Token IDs: vocabulary indices
\item Token log-probabilities: $\log P(y_t\mid\text{context})$
\item Token position: $t \in [1, T]$
\end{itemize}

\section{Distribution Analysis}
% TODO: ~2 pages. Three figures answering three sub-questions:
%   (a) Are token NLLs HETEROGENEOUS within an answer? -> histogram + IQR.
%   (b) Does NLL depend on POSITION? -> mean NLL vs. position t.
%   (c) Does NLL separate CORRECT from INCORRECT answers? -> per-token
%       NLL distribution split by judge label; report t-test / KS.
% Each figure needs: caption, axes, sample size (n), and ONE sentence
% of interpretation. Do NOT foreshadow the fix (weighting) here; just
% describe what you see. The weighting argument lives in \S3.4.
% Analysis on representative model (Llama-3.2-1B or 8B)
Figure 3.1: Token NLL Distribution
> "Token-level NLL values range from 0.1 to 5.2 across all answers, with a mean of 0.8 (SD=0.6). This high variance suggests tokens differ substantially in how confidently the model generates them."
Figure 3.2: NLL by Token Position
> "Early tokens (positions 1-5) show higher NLL (mean=1.2) compared to later tokens (positions 20-25, mean=0.6), indicating the model becomes more confident as generation proceeds."
Figure 3.3: Correct vs Incorrect Answers
> "Correct answers show mean token NLL of 0.7, while incorrect answers show mean NLL of 1.1 (p < 0.001, t-test). This 57 percent increase suggests token-level confidence correlates with answer quality."
\section{Observations}
% TODO: ~0.5 page. A NUMBERED list (F1..F4) of the empirical claims,
% each with an effect size and a figure/table pointer. These bullets
% are the evidence Ch.~7 will cite to answer RQ1. Keep them factual --
% no ``therefore'' or ``this motivates'' clauses; that's \S3.4.
% Summary of the empirical findings only. No re-motivation.
Finding 1: Tokens vary substantially in NLL (range: 0.1-5.2)
Finding 2: Position effects exist (early tokens more uncertain)
Finding 3: Incorrect answers have higher per-token NLL
Finding 4: Not all tokens are equally predictive

\section{Implications for Weighted Aggregation}
% TODO: ~0.25 page. ONE paragraph only. Connect F1-F4 to the next
% chapter -- heterogeneity (F1) + correlation with correctness (F3)
% imply a weighted aggregator can do better than uniform G-NLL, but
% the WEIGHTS must come from a signal the model provides at inference
% time (no ground-truth access). This motivates the ablation-based
% relevance signal formalised in \S\ref{ch:rwgnll}.
% In-chapter bridge into Chapter 4 (NOT a restatement of \S1.1/\S1.2).
Implication: Uniform weighting (G-NLL) may be suboptimal.
A weighting scheme that accounts for token importance could improve
uncertainty estimation.


\chapter{Relevance-Weighted G-NLL (RW-G-NLL)}
\label{ch:rwgnll}
% TODO (chapter scope): ~5-8 pages. This is the main methods chapter.
% Budget split: opener=0.25p, \S4.1=2p, \S4.2=2p, \S4.3=1p, \S4.4=1.5p.
%
% No \section{Problem Formulation} here on purpose: the problem was stated
% in \S1.2 and motivated empirically in Chapter~\ref{ch:token-analysis}.
% Open with 2-3 sentences pointing back to those, then move straight into
% the method. TODO: tighten the paragraph below to a cross-reference.
Given the token-level analysis showing non-uniform importance, we propose Relevance-Weighted G-NLL (RW-G-NLL), which weights each token's log-likelihood by its semantic relevance to the overall answer.
\section{Token Ablation Method}
% TODO: ~2 pages. Define the RELEVANCE SIGNAL formally.
% Cover: (1) how a token is ``ablated'' (mask vs delete vs UNK -- state
% which you use and why); (2) the similarity function sim(y, y_{\setminus t})
% -- embedding model, pooling, cosine vs NLI; (3) how you map similarity
% to a weight w_t (1 - sim, softmax, clipping); (4) edge cases (BOS/EOS,
% whitespace, subword pieces). Include a small worked example (one
% question, 5-10 tokens) with a weight table so the reader SEES the
% mechanism before \S4.2.
\section{The RW-G-NLL Algorithm}
% TODO: ~2 pages. (1) The equation, numbered and referenceable.
% (2) Pseudocode (numpy-style) matching your implementation in the repo.
% (3) Normalisation choice: unnormalised sum vs sum-of-weights
% normalisation vs length-normalised -- state which is the default and
% why. (4) A ``properties'' paragraph: reduces to G-NLL when w_t=const,
% strictly \geq 0, invariant under w_t rescaling IFF normalised.
% RW-G-NLL formula
\section{Computational Considerations}
% TODO: ~1 page. The WHOLE point of RW-G-NLL vs SAR is cost.
% Quantify: SAR needs K stochastic samples + K*T ablation passes
% (~K*(T+1) forwards total); RW-G-NLL needs 1 greedy + T ablation
% forwards. Give a wall-clock table for 1B/8B/70B on your hardware
% (TriviaQA mean answer length). Also discuss: KV-cache reuse for
% ablation, batching strategies, memory (this is useful in Appx.~D too).
\section{Relationship to Existing Methods}
% TODO: ~1.5 pages. Short comparative subsections, not re-explanations:
%   - vs G-NLL: RW-G-NLL reduces to it when w_t = 1/T.
%   - vs length-normalised NLL: same cost, different weighting principle.
%   - vs SAR: same weighting principle, single-sample vs multi-sample.
%   - vs Semantic Entropy: orthogonal (aleatoric token-level vs
%     semantic-cluster epistemic). Note they can be combined.
% End with a small taxonomy figure/table placing RW-G-NLL on a
% (cost, signal type) 2x2.


\chapter{Experimental Design}
\label{ch:exp-design}
% TODO (chapter scope): ~3-5 pages. Reader should be able to replicate
% any number in Ch.~6 after reading this chapter. Do NOT include any
% results or plots here. Budget split: \S5.1=0.75p, \S5.2=0.5p,
% \S5.3=1p, \S5.4=1p, \S5.5=0.5p.
% No \section{Research Hypotheses} here: RQs live in \S1.3 and the
% hypothesis in \S1.4. Open the chapter with one short paragraph of the
% form "We evaluate H1 (\S\ref{...}) through the protocol described
% below." Do NOT restate RQ1--RQ4 or the hypothesis verbatim.
% Models: Small (1B), Medium (8B), Large (70B) - one from each family
\section{Models}
% TODO: ~0.75 page. Focus: justify the model grid, not describe each
% model in depth (that's background). Use 3 scale tiers (e.g. 1B / 8B /
% 70B) and >= 2 families (e.g. Llama-3, Qwen-3) so RQ4 (generalisation)
% has a real answer. Tabulate: family, params, HF id, tokenizer, context
% length. Note quantisation / precision if any (this matters for logprobs).
\section{Datasets}
% TODO: ~0.5 page. Primary: TriviaQA (short factoid, closed-book). Briefly
% mention SQuAD and CoQA as secondary / appendix. For each: split used,
% sample size, why it is a VALID testbed for UQ (i.e. there is a clear
% notion of correctness). Acknowledge: TriviaQA's reference answers are
% imperfect -> motivates LLM-as-judge in \S5.4.
% Dataset: TriviaQA (primary) + brief mention of SQuAD
% Baselines: G-NLL, linear weighting, SAR, Semantic Entropy
\section{Baseline Methods}
% TODO: ~1 page. Three baseline tiers mirroring \S2.4:
%   (B1) Cheap aggregators: G-NLL, length-normalised NLL, PPL.
%   (B2) Positional schemes: linear decay, middle-peak (from your
%        Phase 6 sweep -- pick the best 2-3, document the rest in Appx.~A).
%   (B3) Strong SOTA: SAR (K=10), Semantic Entropy (K=10 + NLI).
% For each baseline, fix ONE set of hyperparameters per paper defaults
% and report those here, not in Ch.~6. State sampling temp for SAR/SE.
% Metrics: AUROC
\section{Evaluation Protocol}
% TODO: ~1 page. Pin down:
%   - Correctness labelling: LLM-as-judge model, prompt template, few-shot
%     examples, tie-breaking rule. Cite your judge-validation run (see
%     \texttt{recompute\_accuracy\_with\_judge.py}).
%   - Metric: AUROC with 95\% CIs via bootstrapping (n=1000 resamples).
%   - Statistical test for pairwise method comparison (paired bootstrap
%     or DeLong's test -- pick one and state it).
%   - Seeds: number of seeds for sampling-based baselines (SAR / SE),
%     deterministic for greedy methods.
\section{Implementation Details}
% TODO: ~0.5 page. Hardware (GPU type and count, see SETUP\_3x11GB\_GPUS.md),
% framework (HF transformers version), precision (bf16/fp16), KV-cache
% reuse for ablation, wandb run groups, random seeds, and a pointer to
% the exact shell scripts (\texttt{run\_small\_models.sh}, \texttt{run\_large\_models.sh},
% \texttt{run\_ultra\_large\_models.sh}). Full config dump goes in Appx.~D.



\chapter{Results and Analysis}
\label{ch:results}
% TODO (chapter scope): ~6-10 pages. One section PER comparison, each
% answering exactly one RQ. Every section should have the same shape:
% (1) headline number / figure, (2) statistical test, (3) 2-3 sentence
% interpretation. Save broader discussion for Ch.~7.
% Budget split: \S6.1=1.5p, \S6.2=1.5p, \S6.3=2p (punchline figure),
% \S6.4=1.5p, \S6.5=1.5p.
\section{RW-G-NLL vs. G-NLL Baseline}
% TODO: ~1.5 pages. Answers RQ2. Main table: AUROC (with 95\% CI) for
% G-NLL, length-normalised NLL, RW-G-NLL per model. Headline figure:
% paired AUROC-delta bar chart. Report win rate over seeds/models. Show
% ONE per-example case where RW-G-NLL flips the ranking correctly.
\section{RW-G-NLL vs. Position-Based Schemes}
% TODO: ~1.5 pages. Answers RQ3. Goal: show that SEMANTIC relevance >
% positional priors. Include the best 2-3 positional schemes from your
% Phase 6 sweep (middle peak, linear decay, NLL-proportional). Full 10+
% schemes go in Appx.~A. Key sub-question: does adding semantic
% relevance on top of a positional prior help? -> one ``hybrid'' row.
% Comparison with best 2-3 positional schemes from Phase 6
\section{RW-G-NLL vs. State-of-the-Art (SAR, Semantic Entropy)}
% TODO: ~2 pages. The efficiency-frontier section. Plot: AUROC vs
% forward-pass count (or wall-clock) across {G-NLL, RW-G-NLL, SAR@K=5,10,
% SE@K=5,10}. The punchline figure of the thesis. State honestly where
% RW-G-NLL loses -- probably to SE on epistemic-heavy queries -- this
% strengthens the ``cheap aleatoric signal'' framing rather than
% weakening it.
\section{Performance Across Model Scales}
% TODO: ~1.5 pages. Answers RQ4. One plot: AUROC vs model size, one
% line per method, one panel per family. Comment on whether the
% RW-G-NLL gain GROWS, SHRINKS, or is FLAT with scale. Acknowledge 70B
% is the largest tested -- generalisation to frontier models is future
% work (Ch.~7).
% Show generalization from 1B to 70B
\section{Weight Pattern Analysis}
% TODO: ~1.5 pages. Qualitative section. Answer ``what tokens does
% RW-G-NLL up-weight?'' by showing (a) part-of-speech distribution of
% top-k weighted tokens, (b) rank correlation between RW-G-NLL weights
% and attention / SAR weights, (c) 2-3 illustrative example answers
% with per-token w_t overlaid. This is where the reader BELIEVES the
% mechanism, so invest one really clean figure here.
% Analysis: what factors make tokens semantically relevant?

\chapter{Discussion and Conclusions}
\label{ch:discussion}
% TODO (chapter scope): ~3-5 pages. This chapter INTERPRETS what
% Ch.~\ref{ch:results} showed; it does NOT introduce new experiments.
% Structure: interpretation -> RQ summary -> limitations -> future
% work -> closing. Every claim here must be backed by a forward/back
% reference to a result section.
% Budget split: \S7.1=1.5p (3 subsections x 0.5p), \S7.2=1p,
% \S7.3=1p, \S7.4=0.5p, \S7.5=0.25p.

\section{Interpretation of Empirical Results}
% TODO: ~1.5 pages total across the three subsections below. Turn
% numbers into narrative. Each subsection takes ONE surprising or
% load-bearing finding and explains WHY it happened, citing the
% relevant result section.
\subsection{The Dominance of Self-Referential Uncertainty}
% Key finding: NLL Proportional (0.724 AUROC) is the most powerful signal, 
% outperforming all other heuristics by identifying "pivotal" error tokens.

\subsection{Positional Priors vs. Semantic Ablation}
% Finding: Middle Peak (0.611 AUROC) outperforming pure SAR (0.596 AUROC) 
% reveals that answer location is a surprisingly robust shortcut.

\subsection{Scaling and Generalization}
% Analysis: Results remain consistent as models scale from 1B to 70B parameters, 
% validating the stability of RW-G-NLL.

\section{Summary of Research Questions}
% TODO: ~1 page. Four short paragraphs (one per RQ). Each is a three-line
% answer: (1) restate RQ as a question, (2) answer YES/NO/PARTIAL with
% the headline number, (3) pointer to the results section. This is where
% RQ1-RQ4 get their definitive answer; ch.~1 posed them, ch.~7 answers.
% RQ1: Establishing token importance heterogeneity.
% RQ2: Confirming AUROC gains via relevance weighting (up to 45% improvement).
% RQ3: Evaluating the synergy of Hybrid Semantic+Positional schemes.
% RQ4: Generalisation across model families and scales (1B-70B).

\section{Practical Implementation and Limitations}
% TODO: ~1 page. Be honest -- reviewers look for this section.
\subsection{The Efficiency Frontier: O(1) vs O(N)}
% Highlighting that NLL Proportional provides SOTA-level UQ without 
% the massive cost of sampling-based methods. 

\subsection{Current Limitations and Threats to Validity}
% TODO: enumerate explicitly. Suggested list:
%   (L1) Requires logit access -> doesn't apply to API-only closed models.
%   (L2) Ablation cost still scales with T -> long-form generation open.
%   (L3) Similarity function is a modelling choice -> sensitivity not
%        exhaustively ablated (forward-ref Appx.~A).
%   (L4) LLM-as-judge correctness labels are noisy; document agreement
%        rate with human labels on a sample.
%   (L5) Evaluated only on English short-form QA; multilinguality open.
% Dependence on internal logprob access and sensitivity to answer length.

\section{Final Contributions and Future Work}
% TODO: ~0.5 page. Contributions = verbatim from \S1.5, now CHECKED-OFF against
% results (one line each: "C1 validated in \S3.4 (F1-F4)", etc.).
% Future work: be concrete, 3-5 bullets. Examples:
%   - RW-G-NLL for abstractive summarisation / code generation.
%   - Combining RW-G-NLL with Semantic Entropy (cheap aleatoric + epistemic).
%   - Learning the weighting function end-to-end via a small head.
%   - Distilling ablation-free approximations of w_t.
% Contribution: A validated methodology for efficient, token-level UQ.
% Future Work: Adapting RW-G-NLL for abstractive summarization.

\section{Concluding Remarks}
% TODO: ~0.25 page. ONE paragraph. Restate the thesis in plain language,
% link back to the motivation in \S1.1, and end on the broader claim
% about cheap, reliable UQ as a prerequisite for safe LLM deployment.
% Closing: Reliable uncertainty is the backbone of safe LLM deployment.

\appendix
\chapter{Additional Weighting Schemes}
% TODO: full grid from your Phase 6 sweep. One table, one figure,
% per-scheme AUROC across models. Text is minimal -- the main chapters
% should cite specific rows by label.
% Detailed results for all 10+ schemes tested
\chapter{Extended Model Results}
% TODO: per-family breakdown (Llama, Qwen, Mistral, ...). Same table
% schema as \S6.4 but exhaustive. Also put prompt-leakage / cleaning
% notes here (see \texttt{fix\_qwen3\_prompt\_leakage.py}).
% Full results for all 9 model families
\chapter{CoQA Dataset Results}
% TODO: secondary-dataset sanity check. Same protocol as TriviaQA, just
% the table. Use this to argue external validity beyond TriviaQA.
\chapter{Implementation Details}
\label{app:impl}
% TODO: reproducibility pack. Repo structure, exact pip env
% (\texttt{nllsar.yml}), hardware, wandb project/entity, seeds, and
% the set of shell scripts needed to regenerate each chapter's figures:
%   Ch.~3 -> \texttt{run\_greedy\_decoding.py}
%   Ch.~6 -> \texttt{run\_multi\_model\_experiments.sh},
%            \texttt{run\_se\_sar\_comparison.sh}
% Code structure, hardware used, and hyperparameter logs

\import{./}{01-introduction}
\import{./}{02-example}
% ...
\import{./}{09-conclusion}

%% 
%%%%%%%%%%%%%%%%%%%%%%%%%%%%%%%%%%%%%%%%%%%%%%%%%%%%%%%%%%%%%%%%%%%%%%%%%%%%%%%%

%%%%%%%%%%%%%%%%%%%%%%%%%%%%%%%%%%%%%%%%%%%%%%%%%%%%%%%%%%%%%%%%%%%%%%%%%%%%%%%%
%% 
%% Print the bibliography
%% 
%% Make sure to start the bibliography on a new odd page (odd is only relevant in twoside layout)
\cleardoubleoddpage
\printbibliography
%% 
%%%%%%%%%%%%%%%%%%%%%%%%%%%%%%%%%%%%%%%%%%%%%%%%%%%%%%%%%%%%%%%%%%%%%%%%%%%%%%%%

%% Begin with the appendix part (all further chapters will be appendices)
\appendix

%%%%%%%%%%%%%%%%%%%%%%%%%%%%%%%%%%%%%%%%%%%%%%%%%%%%%%%%%%%%%%%%%%%%%%%%%%%%%%%%
%% 
%% Include your appendix chapters
%% 

\import{./}{91-appendix}
% ...

%% 
%%%%%%%%%%%%%%%%%%%%%%%%%%%%%%%%%%%%%%%%%%%%%%%%%%%%%%%%%%%%%%%%%%%%%%%%%%%%%%%%

\cleardoubleoddpage

\end{document}
\endinput
